% Dinâmica: 2 blocos num plano inclinado, com ângulos alfa e beta, aceleração e tensão. 
% Faro, seg 15 dez 2025 11:03:29 
% Written by Tordar 
% pstricks2pdf.sh plano_alpha_beta02
\documentclass{article}
\pagestyle{empty}

\usepackage{pst-all} 
\usepackage{pstricks-add}
\usepackage{pst-slpe}
\usepackage{graphicx}
\input{apnum}

% Bloco
\newcommand{\bloco}[5][]{% #1=options, #2=cx, #3=cy, #4=totalWidth, #5=height
  % Left frame (half of total width, full height)
  \psframe[linewidth=1pt,linecolor=gray,fillstyle=slope,slopebegin=white,slopeend=gray,slopesteps=15]
    (!#2 #4 2 div sub #3 #5 2 div sub)                     % lower-left of total area
    (!#2 #3 #5 2 div add)                                  % upper-right of left frame (center x, top)  
  % Right frame (half of total width, full height)
  \psframe[linewidth=1pt,linecolor=gray,fillstyle=slope,slopebegin=gray,slopeend=white,slopesteps=15]
    (!#2 #3 #5 2 div sub)                                  % lower-left of right frame (center x, bottom)
    (!#2 #4 2 div add #3 #5 2 div add)                     % upper-right of total area
  \psframe[#1](!#2 #4 2 div sub #3 #5 2 div sub)(!#2 #4 2 div add #3 #5 2 div add) 
}

\begin{document}

% Altura e largura de cada bloco: 
\evaldef\wb{2}
\evaldef\hb{1}
% Ângulos:
\evaldef\myalpha{60}
\evaldef\mybeta{30}
\evaldef\betai{180-\mybeta}
% Coordenadas do plano inclinado: 
\evaldef\xi{1}
\evaldef\xii{3}
\evaldef\yii{(\xii-\xi)*\TAN{\myalpha*\PI/180}}
\evaldef\xiii{\xii+\yii/(\TAN{\mybeta*\PI/180})}
% Coordenadas para o bloco 1:
\evaldef\rix{(\xii-\xi)/2+\xi}
\evaldef\riy{\yii/2}
\evaldef\dix{-\hb/2*\SIN{\myalpha*\PI/180}}
\evaldef\diy{\hb/2*\COS{\myalpha*\PI/180}}
\evaldef\xbi{\rix+\dix}
\evaldef\ybi{\riy+\diy}
\evaldef\ymi{\hb+0.1}
% Coordenadas para o bloco 2:
\evaldef\diix{\hb/2*\COS{\mybeta*\PI/180}}
\evaldef\diiy{\hb/2*\SIN{\mybeta*\PI/180}}
\evaldef\xbii{\xii+(\xiii-\xii)/2+\diix}
\evaldef\ybii{\yii/2+\diiy}
\evaldef\ymii{\hb+0.1}
% Coordenadas para o fio: 
\evaldef\xint{\xii+\dix+\diy}
\evaldef\yint{\yii+\diiy+\diiy}
% Coordenadas para alpha:
\evaldef\xalpha{\xi+2*\COS{\myalpha*\PI/360}}
\evaldef\yalpha{2*\SIN{\myalpha*\PI/360}}
% Coordenadas para beta:
\evaldef\xbeta{\xiii+2.5*\COS{(180-\mybeta/2)*\PI/180}}
\evaldef\ybeta{2.5*\SIN{\mybeta*\PI/360}}
% Coordenadas para a aceleração: 
\evaldef\ya{\hb+1.5}
\evaldef\yva{\ya-0.5}
\evaldef\xvas{-0.5}
\evaldef\xvae{0.5}
% Coordenadas para a tensão:
\evaldef\xvtmi{0.9*\wb} % bloco 1
\evaldef\xtmi{0.7*\wb}  % bloco 1
\evaldef\ytmi{0.3}      % bloco 1
\evaldef\xvtmii{-0.9*\wb} % bloco 2
\evaldef\xtmii{-0.7*\wb}  % bloco 2
% Coordenadas para os pesos: 
\evaldef\xpi{\xi+(\xii-\xi)/2}
\evaldef\ypi{2+\yii/2}
\evaldef\xpii{\xii+(\xiii-\xii)/2}
% Coordenadas para os ângulos das perpendiculares: 
\evaldef\xalphap{\xpi+0.4}
\evaldef\yalphap{\ypi-0.6}
\evaldef\xbetap{\xpii-0.3}
\evaldef\ybetap{\ypi-1}
%======================================================================
% Let's go:
\begin{pspicture}(0,0)(10,10)
% Chão: 
\psframe[linewidth=0pt,linecolor=lightgray,dimen=inner,fillstyle=hlines*,fillcolor=lightgray,hatchangle=45](0,0)(10,2) 
% Plano inclinado com blocos:
\rput{0}(0,2){ 
%======================================================================
\pspolygon[fillstyle=solid,fillcolor=gray](\xi,0)(\xii,\yii)(\xiii,0) % Plano inclinado 
%\psarc[linewidth=0.8mm,linecolor=black,linestyle=dashed,dash=4mm 1mm]{->}(\xi,0){1.5}{0}{\myalpha} % arc alpha
%\psarc[linewidth=0.8mm,linecolor=black,linestyle=dashed,dash=4mm 1mm]{->}(\xiii,0){2}{\betai}{180} % arc beta
%\rput{0}(\xalpha,\yalpha){\scalebox{1.5}{$\alpha$}}
%\rput{0}(\xbeta,\ybeta){\scalebox{1.5}{$\beta$}}
% Segmentos: 
\psline[linewidth=2pt,linecolor=black](\xii,\yii)(\xint,\yint)
\psline[linewidth=2pt,linecolor=black](\xbi,\ybi)(\xint,\yint)(\xbii,\ybii)
\qdisk(\xint,\yint){2mm}
% Bloco 1:
\rput{\myalpha}(\xbi,\ybi){
\psline[linewidth=3pt,linecolor=blue]{->}(0,-0.05)(\xvtmi,-0.05) % Tensão
\bloco[linewidth=2pt,linecolor=black]{0}{0}{\wb}{\hb}
\rput{0}(0,\ymi){\scalebox{1.5}{$m_1$}}
\rput{0}(0,\ya){\scalebox{1.5}{$\mathbf{a}$}}
\rput{0}(\xtmi,\ytmi){\scalebox{1.5}{$\mathbf{T}$}}
\psline[linewidth=3pt,linecolor=blue]{->}(\xvas,\yva)(\xvae,\yva)
}
% Bloco 2:
\rput{-\mybeta}(\xbii,\ybii){
\psline[linewidth=3pt,linecolor=blue]{->}(0,-0.05)(\xvtmii,-0.05) % Tensão
\bloco[linewidth=2pt,linecolor=black]{0}{0}{\wb}{\hb}
\rput{0}(0,\ymii){\scalebox{1.5}{$m_2$}}
\rput{0}(0,\ya){\scalebox{1.5}{$\mathbf{a}$}}
\rput{0}(\xtmii,\ytmi){\scalebox{1.5}{$\mathbf{T}$}}
\psline[linewidth=3pt,linecolor=blue]{->}(\xvas,\yva)(\xvae,\yva)
}
%======================================================================
}
% Peso 1:
\psline[linewidth=3pt,linecolor=blue]{->}(\xpi,\ypi)(\xpi,2.5)
\rput{\myalpha}(\xpi,\ypi){
\psline[linewidth=0.6mm,linecolor=black,linestyle=dashed,dash=4mm 1mm](0,0)(0,-2) % Perpendicular
} % \rput...
\rput{0}(\xalphap,\yalphap){\scalebox{1}{$\alpha$}}
% Peso 2:
\psline[linewidth=3pt,linecolor=blue]{->}(\xpii,\ypi)(\xpii,2.5)
\rput{-\mybeta}(\xpii,\ypi){
\psline[linewidth=0.6mm,linecolor=black,linestyle=dashed,dash=4mm 1mm](0,0)(0,-1.5) % Perpendicular
} % \rput ... 
\rput{0}(\xbetap,\ybetap){\scalebox{1}{$\beta$}}
\end{pspicture}

\end{document}
